% Phasor hEP Derivation — coupled recurrence formulation
% Compile standalone: pdflatex phasor_hep_derivation.tex

\documentclass[11pt]{article}
\usepackage{amsmath,amssymb,amsthm}
\usepackage{bm}
\usepackage{booktabs}
\usepackage{enumitem}
\usepackage[margin=1in]{geometry}
\usepackage{hyperref}

\newcommand{\R}{\mathbb{R}}
\newcommand{\C}{\mathbb{C}}
\newcommand{\im}{\mathrm{i}}
\newcommand{\dd}{\mathrm{d}}
\newcommand{\dt}{\Delta t}
\newcommand{\eg}{e.g.}
\newcommand{\ie}{i.e.}

\theoremstyle{definition}
\newtheorem{definition}{Definition}
\newtheorem{proposition}{Proposition}
\newtheorem{theorem}{Theorem}
\newtheorem{remark}{Remark}
\newtheorem{corollary}{Corollary}

\DeclareMathOperator{\diag}{diag}
\DeclareMathOperator{\Real}{Re}

\title{Holomorphic Equilibrium Propagation\\
via Coupled Phasor Recurrences}
\author{}
\date{}

\begin{document}
\maketitle

%======================================================================
\section{Equilibrium Propagation: Standard Formulation}\label{sec:ep}
%======================================================================

\subsection{Energy function}

Consider an $L$-layer network with input $\bm{x}\in\C^{d_0}$ (fixed),
hidden states $\bm{z}_1\in\C^{d_1},\dots,\bm{z}_L\in\C^{d_L}$,
weight matrices $W_\ell\in\C^{d_\ell\times d_{\ell-1}}$, biases
$\bm{b}_\ell\in\C^{d_\ell}$, and a holomorphic activation function
$\sigma:\C\to\C$ applied element-wise.  Define $\bm{z}_0\equiv\bm{x}$
and the \emph{pre-activation}
\begin{equation}\label{eq:pre}
  \bm{a}_\ell \;=\; W_\ell\,\bm{z}_{\ell-1} + \bm{b}_\ell,
  \qquad \ell=1,\dots,L.
\end{equation}

The Hopfield-style energy function is
\begin{equation}\label{eq:energy}
  \Phi(\bm{z},\theta,\beta)
  \;=\;
  \sum_{\ell=1}^{L}
    \bigl\langle
      \sigma(\bm{a}_\ell),\;
      \bm{z}_\ell
    \bigr\rangle
  \;-\;
  \beta\,\mathcal{C}(\bm{z}_L,\,\bm{y}),
\end{equation}
where $\theta=\{W_\ell,\bm{b}_\ell\}$ collects all parameters,
$\beta\in\C$ is the nudge strength,
$\mathcal{C}$ is the task cost, $\bm{y}$ is the target, and
$\langle\bm{u},\bm{v}\rangle = \sum_j u_j\,v_j$
is the \textbf{complex bilinear form} (no conjugation).

\begin{remark}
Using the bilinear form rather than the Hermitian inner product
$\bm{u}^H\bm{v}$ is essential:  the Hermitian product involves
complex conjugation and is not holomorphic.  The bilinear form keeps
$\Phi$ holomorphic in all $\bm{z}_\ell$.
\end{remark}

\subsection{Equilibrium condition}\label{sec:eq-cond}

The state $\bm{z}_\ell$ appears in two terms of $\Phi$:
\begin{enumerate}[label=(\alph*)]
\item As the ``post-synaptic'' factor in layer~$\ell$:
  $\langle\sigma(\bm{a}_\ell),\,\bm{z}_\ell\rangle$.
\item As part of the ``pre-synaptic'' input to layer~$\ell+1$
  (when $\ell<L$):
  $\langle\sigma(\bm{a}_{\ell+1}),\,\bm{z}_{\ell+1}\rangle$
  where $\bm{a}_{\ell+1}=W_{\ell+1}\,\bm{z}_\ell+\bm{b}_{\ell+1}$.
\end{enumerate}
Taking the holomorphic gradient $\partial\Phi/\partial\bm{z}_\ell$:

\begin{equation}\label{eq:ep-grad}
  \boxed{
  \frac{\partial\Phi}{\partial\bm{z}_\ell}
  \;=\;
  \underbrace{\sigma(\bm{a}_\ell)}_{\text{drive from below}}
  \;+\;
  \underbrace{
    W_{\ell+1}^{\!\top}\!
    \bigl[
      \sigma'(\bm{a}_{\ell+1})\odot\bm{z}_{\ell+1}
    \bigr]
  }_{\text{drive from above}}
  \;-\;
  \beta\,\delta_{\ell L}\,
  \frac{\partial\mathcal{C}}{\partial\bm{z}_L}
  }
\end{equation}
where $\odot$ is the Hadamard (element-wise) product and
$\delta_{\ell L}$ is the Kronecker delta (the cost gradient acts only
on the output layer).  Note that $W_{\ell+1}^{\!\top}$ (transpose,
not conjugate transpose) couples layer~$\ell+1$ back into
layer~$\ell$ --- this is the \textbf{symmetric weight} requirement of EP.

At equilibrium, $\partial\Phi/\partial\bm{z}_\ell=0$ for all~$\ell$,
and the equilibrium states $\bm{z}_\ell^*(\beta)$ depend on $\beta$.

\subsection{EP gradient theorem}

\begin{theorem}[Scellier \& Bengio, 2017]
If $\bm{z}^*(\beta)$ is a smooth family of equilibria, then
\begin{equation}\label{eq:ep-grad-thm}
  \frac{\dd\mathcal{L}}{\dd\theta}
  \;=\;
  \lim_{\beta\to 0}\;
  \frac{1}{\beta}
  \left[
    \left.\frac{\partial\Phi}{\partial\theta}\right|_{\bm{z}^*(\beta)}
    -
    \left.\frac{\partial\Phi}{\partial\theta}\right|_{\bm{z}^*(0)}
  \right],
\end{equation}
where $\mathcal{L}=\mathcal{C}(\bm{z}_L^*(0),\,\bm{y})$ is the loss
at the free equilibrium.
\end{theorem}

%======================================================================
\section{Phasor Oscillator Dynamics}\label{sec:phasor}
%======================================================================

\subsection{Single-layer ODE}

Each layer's autonomous dynamics follow the resonate-and-fire ODE:
\begin{equation}\label{eq:rf-ode}
  \frac{\dd\bm{z}_\ell}{\dd t}
  \;=\;
  K_\ell\,\bm{z}_\ell + \bm{I}_\ell(t),
  \qquad
  K_\ell = \diag(\bm{k}_\ell),
  \qquad
  k_{\ell,c} = \lambda_{\ell,c} + \im\,\omega_{\ell,c},
\end{equation}
where $\lambda_{\ell,c}<0$ is the per-channel decay rate,
$\omega_{\ell,c}$ is the angular frequency, and $\bm{I}_\ell(t)$ is
the input current.  The diagonal structure means each channel $c$ is an
independent damped oscillator.

\subsection{Discretization}

With time step $\dt$, the exact solution over one step is
\begin{equation}\label{eq:discretize}
  \bm{z}_\ell[n+1]
  \;=\;
  A_\ell\,\bm{z}_\ell[n] + B_\ell\,\bm{I}_\ell[n],
\end{equation}
where $A_\ell = \diag(e^{\bm{k}_\ell\dt})$ and
$B_\ell = \diag\!\bigl((e^{\bm{k}_\ell\dt}-1)/\bm{k}_\ell\bigr)$
are diagonal matrices.  The magnitudes $|A_{\ell,c}|=e^{\lambda_{\ell,c}\dt}<1$,
so the autonomous dynamics are strictly contractive.

\subsection{Impulse-response kernel}

Unrolling the recurrence with zero initial conditions:
\begin{equation}\label{eq:kernel}
  \bm{z}_\ell[n]
  \;=\;
  \sum_{j=0}^{n} \mathcal{K}_\ell[n-j]\,\bm{I}_\ell[j],
  \qquad
  \mathcal{K}_\ell[m] = A_\ell^m\,B_\ell.
\end{equation}
This is a causal convolution with the phasor kernel
$\mathcal{K}_\ell[m]$, computable via FFT.  The kernel is a decaying
spiral in the complex plane---each channel $c$ traces
$\mathcal{K}_{\ell,c}[m] = e^{k_{\ell,c}\,m\dt}\,(e^{k_{\ell,c}\dt}-1)/k_{\ell,c}$.

%======================================================================
\section{Coupled Phasor Recurrence for EP}\label{sec:coupled}
%======================================================================

\subsection{Merging EP dynamics with phasor settling}

We now replace EP's generic fixed-point iteration with the phasor
oscillator dynamics.  The ``input current'' to each layer's oscillator
bank is the EP energy gradient~\eqref{eq:ep-grad}:

\begin{equation}\label{eq:coupled-drive}
  \bm{I}_\ell[n]
  \;=\;
  \underbrace{\sigma(\bm{a}_\ell[n])}_{\text{feedforward}}
  \;+\;
  \underbrace{
    W_{\ell+1}^{\!\top}
    \bigl[\sigma'(\bm{a}_{\ell+1}[n])\odot\bm{z}_{\ell+1}[n]\bigr]
  }_{\text{feedback}}
  \;-\;
  \beta\,\delta_{\ell L}\,
  \frac{\partial\mathcal{C}}{\partial\bm{z}_L}\bigg|_{\bm{z}[n]}
\end{equation}
where $\bm{a}_\ell[n]=W_\ell\,\bm{z}_{\ell-1}[n]+\bm{b}_\ell$ is
the pre-activation computed from the current states.

The coupled recurrence is then:

\begin{equation}\label{eq:coupled}
  \boxed{
  \bm{z}_\ell[n+1]
  \;=\;
  A_\ell\,\bm{z}_\ell[n]
  \;+\;
  B_\ell\,\bm{I}_\ell[n],
  \qquad \ell=1,\dots,L.
  }
\end{equation}

All layers update \textbf{simultaneously} at each time step~$n$.
The oscillator eigenvalues $\bm{k}_\ell$ provide the
settling dynamics:
\begin{itemize}
\item The \textbf{decay} ($\lambda<0$) drives the system toward
  equilibrium---it acts as a natural damping that replaces the step-size
  parameter $\eta$ in standard EP.
\item The \textbf{oscillation} ($\omega$) adds phasor structure---phase
  information is carried implicitly in the complex state without
  explicit phase extraction.
\item The \textbf{kernel} $\mathcal{K}_\ell[m]=A_\ell^m B_\ell$
  integrates the EP gradient signal over time with exponential memory,
  smoothing the settling trajectory.
\end{itemize}

\subsection{Two-layer example}

For concreteness, consider $L=2$ (hidden layer~$\bm{z}_1$, output
layer~$\bm{z}_2$, input $\bm{x}=\bm{z}_0$ fixed):

\begin{align}
  \bm{I}_1[n] &= \sigma(W_1\bm{x}+\bm{b}_1)
    + W_2^{\!\top}\bigl[\sigma'(W_2\bm{z}_1[n]+\bm{b}_2)
      \odot\bm{z}_2[n]\bigr],
    \label{eq:I1}\\[4pt]
  \bm{I}_2[n] &= \sigma(W_2\bm{z}_1[n]+\bm{b}_2)
    - \beta\,\frac{\partial\mathcal{C}}{\partial\bm{z}_2}\bigg|_{\bm{z}[n]},
    \label{eq:I2}\\[6pt]
  \bm{z}_1[n+1] &= A_1\,\bm{z}_1[n] + B_1\,\bm{I}_1[n],
    \label{eq:z1}\\
  \bm{z}_2[n+1] &= A_2\,\bm{z}_2[n] + B_2\,\bm{I}_2[n].
    \label{eq:z2}
\end{align}

Note the symmetric coupling:  $W_2$ sends information forward from
$\bm{z}_1$ to $\bm{z}_2$ in~\eqref{eq:I2}, and $W_2^{\!\top}$ sends
influence backward from $\bm{z}_2$ to $\bm{z}_1$ in~\eqref{eq:I1}.

\subsection{Equilibrium}

At a fixed point ($\bm{z}_\ell[n+1]=\bm{z}_\ell[n]=\bm{z}_\ell^*$):
\begin{equation}\label{eq:fixed-point}
  \bm{z}_\ell^*
  \;=\;
  A_\ell\,\bm{z}_\ell^*+B_\ell\,\bm{I}_\ell^*
  \qquad\Longrightarrow\qquad
  \bm{z}_\ell^*
  \;=\;
  (I-A_\ell)^{-1}\,B_\ell\,\bm{I}_\ell^*.
\end{equation}
Since $|A_{\ell,c}|<1$, the matrix $(I-A_\ell)$ is invertible with
$(I-A_\ell)^{-1}_{cc} = 1/(1-e^{k_{\ell,c}\dt})$.  The fixed-point
condition~\eqref{eq:fixed-point} is a nonlinear system (because
$\bm{I}_\ell^*$ depends on all $\bm{z}^*$), solved self-consistently
by running the recurrence until convergence.

\begin{remark}[Static gain]
The factor $(I-A_\ell)^{-1}B_\ell$ is the \emph{DC gain} of the
oscillator---the steady-state response to a constant input.
Channel~$c$ has DC gain
\[
  G_c = \frac{B_c}{1-A_c}
  = \frac{e^{k_c\dt}-1}{k_c(1-e^{k_c\dt})}
  = -\frac{1}{k_c}
  = \frac{1}{|\lambda_c|^2+\omega_c^2}
    \bigl(|\lambda_c|-\im\omega_c\bigr).
\]
The magnitude $|G_c|=1/|k_c|$ is larger for slowly-decaying,
low-frequency channels (small $|\lambda|$, small $\omega$),
which act as long-memory integrators.
\end{remark}

\subsection{Convergence}

\begin{proposition}[Contraction]
If $\sigma$ has bounded derivative ($|\sigma'|\le M$) and the coupling
is weak enough that
\begin{equation}\label{eq:contraction}
  \max_\ell\;|A_{\ell,c}| + |B_{\ell,c}|\,M\,\|W\|_{\text{op}}
  \;<\; 1
  \qquad\text{for all channels $c$},
\end{equation}
then the coupled recurrence~\eqref{eq:coupled} is a contraction
mapping and converges to a unique fixed point.
\end{proposition}

In practice, the damping rates $|\lambda_{\ell,c}|$ and the weight
magnitudes determine whether the system converges.  The holomorphic
activation $\sigma=\tanh$ has $|\tanh'|\le 1$, so the contraction
condition reduces to $|A_{\ell,c}|+|B_{\ell,c}|\,\|W\|_{\text{op}}<1$.

%======================================================================
\section{Continuous-Time Formulation}\label{sec:ode}
%======================================================================

\subsection{Coupled ODE system}

The discrete recurrence~\eqref{eq:coupled} corresponds to the
continuous-time coupled ODE system:
\begin{equation}\label{eq:coupled-ode}
  \boxed{
  \frac{\dd\bm{z}_\ell}{\dd t}
  \;=\;
  K_\ell\,\bm{z}_\ell
  \;+\;
  \bm{I}_\ell(t),
  \qquad \ell=1,\dots,L,
  }
\end{equation}
where $K_\ell=\diag(\bm{k}_\ell)$ and $\bm{I}_\ell(t)$ is the EP
energy gradient~\eqref{eq:coupled-drive} evaluated at the instantaneous
states.  This is a system of $\sum_\ell d_\ell$ coupled complex ODEs.

The connection to the existing \texttt{oscillator\_bank} is direct:
the oscillator bank solves $\dd\bm{z}/\dd t = K\bm{z}+\bm{I}(t)$
for a single layer.  The coupled formulation simply makes
$\bm{I}_\ell(t)$ depend on all layers' states, creating an interacting
system.

\subsection{Lyapunov stability}

For the autonomous system ($\bm{I}=0$), the Lyapunov function is
$V = \sum_\ell \|\bm{z}_\ell\|^2$, with
$\dot{V} = 2\sum_\ell\Real(\bm{z}_\ell^H K_\ell\bm{z}_\ell)
= 2\sum_{\ell,c}\lambda_{\ell,c}|z_{\ell,c}|^2 < 0$
since all $\lambda_{\ell,c}<0$.  The autonomous dynamics are globally
asymptotically stable.

When the coupling $\bm{I}_\ell$ is present, local stability is
governed by the spectrum of the full Jacobian:
\[
  J = \begin{pmatrix}
    K_1 & \partial\bm{I}_1/\partial\bm{z}_2 \\
    \partial\bm{I}_2/\partial\bm{z}_1 & K_2
  \end{pmatrix}
\]
(for $L=2$).  Stability requires all eigenvalues of $J$ to have
negative real part, which holds when the damping dominates the
coupling strength---the same condition as~\eqref{eq:contraction}
in continuous time.

%======================================================================
\section{Holomorphic Extension}\label{sec:holo}
%======================================================================

\subsection{Complex nudge parameter}

Following Laborieux \& Zenke~(2022), we extend $\beta$ from
$\R$ to $\C$.  If $\Phi$ is holomorphic in $\beta$ and
the equilibrium map $\bm{z}^*(\beta)$ is holomorphic, then
by the Cauchy integral formula:
\begin{equation}\label{eq:cauchy}
  \left.\frac{\dd}{\dd\beta}
    \frac{\partial\Phi}{\partial\theta}\right|_{\beta=0}
  \;=\;
  \frac{1}{2\pi\im}
  \oint_{|\beta|=r}
    \frac{1}{\beta^2}\,
    \frac{\partial\Phi}{\partial\theta}\bigg|_{\bm{z}^*(\beta)}
    \dd\beta.
\end{equation}

Discretizing the contour with $N$ equally-spaced points
$\beta_n = r\,e^{2\pi\im n/N}$, $n=0,\dots,N-1$:
\begin{equation}\label{eq:fourier}
  \boxed{
  \frac{\dd\mathcal{L}}{\dd\theta}
  \;=\;
  \frac{1}{N}\sum_{n=0}^{N-1}
    \left.\frac{\partial\Phi}{\partial\theta}\right|_{\bm{z}^*(\beta_n)}
    \!\!\cdot\, e^{-2\pi\im n/N}.
  }
\end{equation}

This is exact for any $r>0$, provided holomorphicity holds inside the
disk $|\beta|\le r$.

\subsection{Holomorphicity requirements}\label{sec:holo-req}

For the energy~\eqref{eq:energy} to be holomorphic in $\bm{z}$ and
$\beta$, we need:

\begin{enumerate}
\item \textbf{Activation $\sigma$}: Must be holomorphic (complex-differentiable).
  Complex $\tanh(z) = (e^z-e^{-z})/(e^z+e^{-z})$ is meromorphic
  (holomorphic except at isolated poles on the imaginary axis at
  $z=\im\pi(n+\tfrac{1}{2})$).  For bounded states, this is sufficient.

\item \textbf{Bilinear form}: The inner product
  $\langle\bm{u},\bm{v}\rangle=\sum_j u_j v_j$ is holomorphic.
  The Hermitian form $\bm{u}^H\bm{v}=\sum_j\bar{u}_j v_j$ is
  \textbf{not} holomorphic due to conjugation.

\item \textbf{Weight operations}: Matrix multiplication
  $W\bm{z}$ is holomorphic in $\bm{z}$ (and in $W$ if weights are
  viewed as variables).  The transpose $W^{\!\top}$ (not $W^H$) is
  used in the feedback term.

\item \textbf{Oscillator dynamics}: The recurrence
  $\bm{z}[n+1]=A\bm{z}[n]+B\bm{I}[n]$ is holomorphic in
  $\bm{z}$ and $\bm{I}$ (linear).  It is also holomorphic in the
  eigenvalues $\bm{k}$ through $A=e^{\bm{k}\dt}$ and
  $B=(e^{\bm{k}\dt}-1)/\bm{k}$ (entire functions of $\bm{k}$).

\item \textbf{Cost function $\mathcal{C}$}: Should be holomorphic
  in $\bm{z}_L$.  See Section~\ref{sec:cost}.
\end{enumerate}

The standard phasor activation \texttt{normalize\_to\_unit\_circle},
$z\mapsto z/|z|$, is \textbf{not} holomorphic because
$|z|=\sqrt{z\bar{z}}$ involves conjugation.  In the coupled
recurrence formulation, this activation is \textbf{not applied during
settling}---the oscillator dynamics and holomorphic $\sigma$ handle
boundedness.  Phase information passes implicitly through the complex
state; explicit phase extraction occurs only at the final readout.

\subsection{Holomorphic cost function}\label{sec:cost}

The cost function enters the energy as $-\beta\,\mathcal{C}(\bm{z}_L,\bm{y})$.
For $\Phi$ to be holomorphic in $\bm{z}_L$, we need $\mathcal{C}$ to
be holomorphic.  Several options:

\begin{enumerate}
\item \textbf{Complex linear cost}:
  $\mathcal{C} = -\langle\bm{t},\,\bm{z}_L\rangle
  = -\sum_j t_j\,z_{L,j}$,
  where $\bm{t}\in\C^{d_L}$ is a target vector (e.g., class-specific
  codebook entry on the unit circle).
  This is holomorphic and linear---the simplest option.

\item \textbf{Softmax cross-entropy on linear readout}:
  Introduce a readout matrix $W_r\in\R^{C\times d_L}$ and define
  logits $\bm{\ell}=W_r\bm{z}_L$ (real matrix times complex vector
  gives complex logits; take real part for classification).
  The softmax and cross-entropy applied to $\Real(\bm{\ell})$ are not
  holomorphic, but applied to the complex logits directly:
  \[
    \mathcal{C} = -\sum_c y_c\,\log\frac{e^{\ell_c}}{\sum_{c'}e^{\ell_{c'}}}
  \]
  is holomorphic in $\bm{\ell}$ (and hence in $\bm{z}_L$), since
  $\exp$ and $\log$ are holomorphic.

\item \textbf{Quadratic cost}:
  $\mathcal{C} = \tfrac{1}{2}\langle\bm{z}_L-\bm{t},\,
  \bm{z}_L-\bm{t}\rangle
  = \tfrac{1}{2}\sum_j(z_{L,j}-t_j)^2$.
  This is holomorphic (note: $(z-t)^2$ is a polynomial in $z$, not
  $|z-t|^2$ which would require conjugation).  However, for complex $z$,
  this does not measure Euclidean distance---it is a complex-valued
  ``distance'' that can be negative.
\end{enumerate}

Option~2 (complex softmax cross-entropy) is the closest to standard
practice and was used in the reference implementation.

%======================================================================
\section{Teaching Oscillation Interpretation}\label{sec:osc}
%======================================================================

\subsection{Oscillating nudge}

In the continuous-time formulation~\eqref{eq:coupled-ode}, let $\beta$
oscillate at a \emph{teaching frequency} $\omega_T$:
\begin{equation}\label{eq:beta-osc}
  \beta(t) = r\,e^{\im\omega_T t}.
\end{equation}

The output layer's ODE becomes
\begin{equation}
  \frac{\dd\bm{z}_L}{\dd t}
  = K_L\bm{z}_L
  + \sigma(W_L\bm{z}_{L-1}+\bm{b}_L)
  - r\,e^{\im\omega_T t}\,
    \frac{\partial\mathcal{C}}{\partial\bm{z}_L}.
\end{equation}
This is a periodically driven oscillator.  After transients decay
(governed by $\lambda<0$), every state $\bm{z}_\ell(t)$ reaches a
periodic steady state that can be expanded in a Fourier series in
$\omega_T$:
\begin{equation}\label{eq:fourier-state}
  \bm{z}_\ell(t)
  \;=\;
  \bm{c}_{\ell,0}
  + \bm{c}_{\ell,1}\,e^{\im\omega_T t}
  + \bm{c}_{\ell,2}\,e^{2\im\omega_T t}
  + \cdots
\end{equation}

\subsection{Gradient as first Fourier coefficient}

The parameter gradient~\eqref{eq:fourier} corresponds to extracting
the first Fourier coefficient of
$\partial\Phi/\partial\theta$ evaluated along the oscillating trajectory.
In continuous time:
\begin{equation}\label{eq:grad-fourier}
  \frac{\dd\mathcal{L}}{\dd\theta}
  \;=\;
  \frac{1}{T}\int_0^T
    \left.\frac{\partial\Phi}{\partial\theta}\right|_{\bm{z}(t)}
    e^{-\im\omega_T t}\,\dd t,
  \qquad T=\frac{2\pi}{\omega_T}.
\end{equation}

This is a \textbf{lock-in measurement}: correlate the energy gradient
with a reference signal at the teaching frequency.  All other frequency
components (DC, harmonics) average out.

\subsection{Physical interpretation}

Each neuron's potential $\bm{z}_\ell(t)$ traces a trajectory in the
complex plane.  The zeroth-order term $\bm{c}_{\ell,0}$ is the free
equilibrium (the prediction).  The first-order term
$\bm{c}_{\ell,1}\,e^{\im\omega_T t}$ is an ``epicycle''---a secondary
oscillation at the teaching frequency, superimposed on the natural
dynamics.

\begin{itemize}
\item The \textbf{amplitude} $|\bm{c}_{\ell,1}|$ encodes the gradient
  magnitude---how sensitive this neuron is to the output error signal.
\item The \textbf{phase} $\arg(\bm{c}_{\ell,1})$ encodes the gradient
  direction---whether the weight should increase or decrease.
\end{itemize}

For the linear oscillator dynamics $\dd z/\dd t = kz + I(t)$, the
response at $\omega_T$ is $c_1 = I_1/(\im\omega_T - k)$, where $I_1$
is the first Fourier coefficient of the input.  The linear transfer
function ensures that the teaching signal at $\omega_T$ and the natural
dynamics at $\omega$ do not interfere: they are orthogonal spectral
components.  This \textbf{spectral separation} is exact in linear
layers and approximate (with controlled error) through nonlinear
activations.

%======================================================================
\section{Parameter Gradients}\label{sec:param-grad}
%======================================================================

\subsection{Energy gradient w.r.t.\ weights}

From~\eqref{eq:energy}, the gradient of $\Phi$ w.r.t.\ $W_\ell$ is:
\begin{equation}\label{eq:dPhi-dW}
  \frac{\partial\Phi}{\partial W_\ell}
  \;=\;
  \bigl[\sigma'(\bm{a}_\ell)\odot\bm{z}_\ell\bigr]
  \,\bm{z}_{\ell-1}^{\!\top}.
\end{equation}
This is a \textbf{Hebbian} outer product: the gradient of the energy
w.r.t.\ a weight is the product of post-synaptic activity
(modulated by the activation derivative) and pre-synaptic activity.
It is local---each synapse needs only the activities of the two neurons
it connects.

\subsection{Energy gradient w.r.t.\ oscillator parameters}

The energy depends on $\bm{k}_\ell=\bm{\lambda}_\ell+\im\bm{\omega}_\ell$
through the equilibrium states $\bm{z}^*(\bm{k})$.  Since the kernel
$\mathcal{K}[m]=A^m B$ is holomorphic in $\bm{k}$, the equilibrium
map inherits this holomorphicity, and we can compute
$\partial\mathcal{L}/\partial\bm{k}_\ell$ via the same contour
integration~\eqref{eq:fourier}.

Equivalently, in the continuous-time formulation, the ODE
$\dd z_c/\dd t = k_c z_c + I_c(t)$ has solution
\[
  z_c(t) = e^{k_c t}z_c(0) + \int_0^t e^{k_c(t-s)}I_c(s)\,\dd s.
\]
This is entire in $k_c$, so holomorphic EP can train the dynamics
parameters $(\lambda_c,\omega_c)$ alongside the weights.

\subsection{Full hEP parameter update}

The complete training step:
\begin{enumerate}
\item \textbf{Free phase}: Run coupled recurrence~\eqref{eq:coupled}
  with $\beta=0$ to convergence.  Record equilibrium
  $\bm{z}^*(0)$.
\item \textbf{Contour evaluation}: For $n=0,\dots,N-1$:
  \begin{enumerate}[label=(\alph*)]
    \item Set $\beta_n = r\,e^{2\pi\im n/N}$.
    \item Initialize from $\bm{z}^*(0)$ (complexified).
    \item Run coupled recurrence to convergence at $\beta_n$.
    \item Compute
      $\bm{g}_n = \partial\Phi/\partial\theta|_{\bm{z}^*(\beta_n)}$.
  \end{enumerate}
\item \textbf{Extract gradient}:
  $\Delta\theta = \frac{1}{N}\sum_n \Real\!\bigl(\bm{g}_n\,
  e^{-2\pi\im n/N}\bigr)$.
\item \textbf{Update}: $\theta \leftarrow \theta - \eta\,\Delta\theta$.
\end{enumerate}

%======================================================================
\section{Connection to PhasorNetworks Dispatch Paths}\label{sec:dispatch}
%======================================================================

\subsection{Mapping to 3D complex dispatch}

The coupled recurrence~\eqref{eq:coupled} maps naturally onto the
existing 3D complex dispatch path in PhasorNetworks.jl.  In the
standard (feedforward) case:
\begin{center}
\begin{tabular}{lll}
\toprule
\textbf{Standard SSM} & \textbf{Coupled hEP} \\
\midrule
Input: $(C_{\text{in}}, L, B)$ complex & Input: $(C_{\text{in}}, B)$,
  $L$ = settling steps \\
$H = W \cdot x$ (channel mixing) & $\bm{I}_\ell[n]$ from
  Eq.~\eqref{eq:coupled-drive} \\
$Z = \text{causal\_conv}(K, H)$ & $\bm{z}_\ell$ from
  recurrence~\eqref{eq:coupled} \\
$Y = \text{activation}(Z)$ & Equilibrium $\bm{z}^*$ (no activation
  during settling) \\
\bottomrule
\end{tabular}
\end{center}

The key difference: in the standard SSM, each layer processes a
pre-computed input sequence.  In the coupled hEP formulation, the
``input'' at each step depends on the current states of all layers,
creating a recurrent interaction.

\subsection{Activation and phase extraction}

In the standard dispatch, the activation
(\texttt{normalize\_to\_unit\_circle}) is applied at each layer's
output, and \texttt{complex\_to\_angle} extracts phase for the Phase
dispatch.

In the hEP formulation:
\begin{itemize}
\item During settling: \textbf{no activation applied}.  The oscillator
  damping ($\lambda<0$) keeps states bounded.  The holomorphic
  activation $\sigma$ (e.g., $\tanh$) appears only inside the energy
  gradient~\eqref{eq:coupled-drive}.
\item At readout: phase can be extracted from the equilibrium states
  $\bm{z}^*$ for compatibility with the Codebook/SSMReadout layers.
  This non-holomorphic step occurs \textbf{outside} the EP dynamics
  and does not affect gradient quality.
\end{itemize}

\subsection{Connection to the continuous ODE dispatch}

The coupled ODE~\eqref{eq:coupled-ode} maps directly onto the
\texttt{oscillator\_bank} infrastructure.  Each layer already solves
$\dd\bm{z}/\dd t = K\bm{z} + \bm{I}(t)$ via
\texttt{DifferentialEquations.jl}.  The hEP extension requires:
\begin{enumerate}
\item Making $\bm{I}_\ell(t)$ depend on other layers' states
  (creating a coupled system).
\item Adding the teaching term
  $-\beta(t)\,\partial\mathcal{C}/\partial\bm{z}_L$ to the output
  layer.
\end{enumerate}
These modifications preserve the ODE structure and remain compatible
with the existing solver and sensitivity analysis infrastructure.

%======================================================================
\section{Consistency Checks}\label{sec:checks}
%======================================================================

\subsection{Recovery of standard EP}

Setting $K_\ell=\lambda\,I$ (uniform real eigenvalues, no oscillation)
and $\dt=1$, the recurrence becomes
\[
  \bm{z}_\ell[n+1]
  = e^{\lambda}\bm{z}_\ell[n]
  + \frac{e^\lambda-1}{\lambda}\,\bm{I}_\ell[n].
\]
As $\lambda\to 0^-$, we have $e^\lambda\to 1$ and
$(e^\lambda-1)/\lambda\to 1$, so
$\bm{z}_\ell[n+1]\approx\bm{z}_\ell[n]+\bm{I}_\ell[n]$---this is
standard gradient ascent on the energy with step size~1.
The phasor formulation thus reduces to standard EP when the oscillator
dynamics are removed.

\subsection{Recovery of feedforward SSM}

If the feedback terms are removed
($W_{\ell+1}^{\!\top}[\cdots]=0$ in~\eqref{eq:coupled-drive}) and
$\beta=0$, then each layer receives only feedforward input:
$\bm{I}_\ell[n]=\sigma(W_\ell\bm{z}_{\ell-1}[n]+\bm{b}_\ell)$.
This is the standard SSM causal convolution, where each layer
processes the previous layer's output as a time series.

\subsection{Holomorphicity of the full system}

Each component of the coupled recurrence is holomorphic:
\begin{itemize}
\item $A_\ell\bm{z}_\ell$: linear in $\bm{z}_\ell$ and
  holomorphic in $\bm{k}_\ell$ (via $A=e^{k\dt}$).
\item $B_\ell\bm{I}_\ell$: $B_\ell$ is holomorphic in $\bm{k}_\ell$;
  $\bm{I}_\ell$ is holomorphic in $\bm{z}$ if $\sigma$ and
  $\mathcal{C}$ are holomorphic.
\item Weight operations: $W\bm{z}$ is holomorphic; $W^{\!\top}\bm{v}$
  is holomorphic (transpose, not conjugate transpose).
\item $\beta$: enters linearly in the cost term.
\end{itemize}
By composition, the equilibrium map
$\bm{z}^*(\beta,\theta)$ is holomorphic in $\beta$ (given $\sigma$
and $\mathcal{C}$ are holomorphic), validating the contour
integration~\eqref{eq:fourier}.

\subsection{Dimensions and types}

\begin{center}
\begin{tabular}{llll}
\toprule
\textbf{Quantity} & \textbf{Type} & \textbf{Shape} & \textbf{Notes} \\
\midrule
$\bm{z}_\ell$ & $\C^{d_\ell}$ & $(d_\ell, B)$ & Hidden states (batch $B$) \\
$W_\ell$ & $\R^{d_\ell\times d_{\ell-1}}$ & --- & Real weights \\
$\bm{b}_\ell$ & $\C^{d_\ell}$ & $(d_\ell,)$ & Complex bias \\
$K_\ell$ & $\C^{d_\ell\times d_\ell}$ & diagonal & Per-channel eigenvalues \\
$A_\ell$ & $\C^{d_\ell\times d_\ell}$ & diagonal & $\diag(e^{\bm{k}_\ell\dt})$ \\
$B_\ell$ & $\C^{d_\ell\times d_\ell}$ & diagonal & $\diag((e^{\bm{k}_\ell\dt}-1)/\bm{k}_\ell)$ \\
$\beta$ & $\C$ & scalar & Nudge (real for std EP) \\
\bottomrule
\end{tabular}
\end{center}

Weights $W_\ell$ are real-valued but operate on complex states,
producing complex outputs.  This matches the existing PhasorDense
convention where \texttt{params.weight} is \texttt{Float32} and the
complex multiplication is done component-wise:
$W\bm{z} = W\,\Real(\bm{z}) + \im\,W\,\text{Im}(\bm{z})$.

%======================================================================
\section{Summary}\label{sec:summary}
%======================================================================

The coupled phasor recurrence~\eqref{eq:coupled} unifies equilibrium
propagation with the phasor SSM framework:

\begin{enumerate}
\item EP's settling dynamics are realized by \textbf{damped oscillator
  recurrences}, with the oscillator eigenvalue $k=\lambda+\im\omega$
  providing both convergence (decay) and phasor structure (oscillation).

\item The EP energy gradient drives each oscillator as an
  \textbf{input current}, with feedforward terms ($\sigma(W\bm{z}_{below})$)
  and feedback terms ($W^{\!\top}[\sigma'\odot\bm{z}_{above}]$)
  coupling all layers simultaneously.

\item The \textbf{holomorphic extension} (complex $\beta$) is naturally
  compatible because the oscillator dynamics, weight operations, and
  holomorphic activations are all complex-differentiable.  The gradient
  is the first Fourier coefficient of the energy gradient evaluated
  along a contour in the complex $\beta$-plane.

\item In continuous time, the oscillating nudge
  $\beta(t)=r\,e^{\im\omega_T t}$ produces an \textbf{epicyclic
  perturbation} of each neuron's trajectory, whose amplitude and phase
  at $\omega_T$ encode the parameter gradient---extractable via lock-in
  measurement.

\item The formulation \textbf{reduces to} standard EP (when
  $\omega=0$, $\lambda\to 0$) and to the feedforward SSM (when feedback
  is removed and $\beta=0$), confirming consistency with both limiting
  cases.
\end{enumerate}

\section*{References}

\begin{itemize}
\item Laborieux, A.\ \& Zenke, F.\ (2022).
  ``Holomorphic Equilibrium Propagation Computes Exact Gradients
  Through Finite Size Oscillations.'' NeurIPS 2022.
  arXiv:2209.00530.
\item Scellier, B.\ \& Bengio, Y.\ (2017).
  ``Equilibrium Propagation: Bridging the Gap Between Energy-Based
  Models and Backpropagation.''
  Frontiers in Computational Neuroscience.
\end{itemize}

\end{document}
